\chapter{Appendix}


\section{Proof of \ref{it: df} \texorpdfstring{\(\Rightarrow\)}{⇒} \ref{it: s.open} from the Portmanteau Theorem}
\label{sec: appendix portmanteau}

We require four facts that would be too much work to develop here. Especially Fatou's lemma
which is a fundamental result in measure theory.

\begin{fact}[Fatou's Lemma]
    \label{fact: fatou}
    Let \((f_n)_{n\in \nat}\) be a sequence of non-negative measurable functions.
    Then for any measure \(\mu\) it holds that
    \[
        \int \liminf_{n\to\infty} f_n \, d\mu
        \le \liminf_{n\to\infty} \int f_n \, d\mu.
    \]
\end{fact}
\begin{fact}
    \label{fact: non-decreasing -> discontinuities countable}
    Let \(f\) be a non-decreasing function. Then its set of discontinuities 
    is at most countable.
\end{fact}
\begin{proof}[Proofsketch]
    The only discontinuities of a non-decreasing function are jumps.
    On bounded intervals there are a finite number of jumps greater than
    \(1/n\). Consequently, there are only countably many jumps of any size
    on such a bounded interval. Since the real line may be covered by a
    countable number of bounded intervals the claim follows.
\end{proof}
\begin{fact}
    \label{fact: discontinuity are atoms}
    \(\Pr(X =x) > 0\) if an only if \(x\) is a discontinuity point of the
    distribution function \(F(x) = \Pr(X \le x)\).
\end{fact}
\begin{proof}[Proofsketch]
    Use the continuity of measure to determine the value of the left-limits 
    of \(F\).
\end{proof}
\begin{fact}
    \label{fact: open set as union of intervals}
    Any open set \(O\subseteq \real\) may be written as a countable,
    disjoint union of open intervals \(O = \biguplus_{i=1}^\infty (a_i,
    b_i)\) for \(a_i, b_i \in \real\).\footnote{
        \(b_i \le a_i\) implies \((a_i, b_i) = \emptyset\).
    }
\end{fact}

Using these facts the proof has two steps: First we prove it for all open
intervals and then use that any open set \(O\) may be written as a
countable, disjoint union of open intervals
\(O = \biguplus_{i=1}^\infty (a_i, b_i)\) for
\(a_i, b_i \in \real\) (Fact \ref{fact: open set as union of intervals}). Thus
by Fatou's lemma (Fact \ref{fact: fatou}) we have
\[\begin{aligned}
    \liminf_{n\to \infty} \Prob{X_n \in O}
    &= \liminf_{n\to \infty} \sum_{i=1}^\infty \Prob[\big]{X_n \in (a_i, b_i)}
    \\
    \overset{\text{Fatou}}&\ge \sum_{i=1}^\infty \liminf_{n\to \infty} \Prob[\big]{X_n \in (a_i, b_i)}
    \\
    \overset{(*)}&\ge \sum_{i=1}^\infty \Prob{X \in (a_i, b_i)}
    &= \Prob{X \in O},
\end{aligned}\]
Here \((*)\) is the claim for open intervals. 
To prove it for an open interval \((a,b)\) pick two sequences \((a_k)_{k\in \nat}\) and \((b_k)_{k\in \nat}\)
of \underline{continuity} points of the distribution function \(F(x) \coloneq \Prob{X \le x}\)
with \(a_k \downarrow a\) and \(b_k \uparrow b\). This is possible since 
\(F\) is non-decreasing and Fact \ref{fact: non-decreasing -> discontinuities countable}.
By \ref{it: df} and Fact \ref{fact:
discontinuity are atoms} we thus have for \(F_n(x) \coloneq \Prob{X_n \le x}\)
that \(F_n(a_k) \to F(a_k)\)
as the \(a_k\) are continuity points and similarly for \(b_k\).
Consequently,
as \((a,b) \supseteq (a_k, b_k]\) we have
\[\begin{aligned}
    \liminf_{n\to \infty} \Prob{X_n \in (a,b)}
    &\ge \liminf_{n\to\infty} \Prob{X_n \in (a_k, b_k]}
    \\
    &= \liminf_{n\to\infty} F_n(b_k) - F_n(a_k)
    \\
    \overset{\ref{it: df}}&= F(b_k) - F(a_k)
    = \Prob{X \in (a_k, b_k]}.
\end{aligned}\]
Since it is true for any \(k\) we have
\[\begin{aligned}
    \liminf_{n\to \infty} \Prob{X_n \in (a,b)}
    &\ge \lim_{k\to\infty} \Prob{X \in (a_k, b_k]}
    \\
    &= \Prob[\Big]{X \in \bigcup_{k\in \nat} (a_k, b_k]}
    = \Prob{X \in (a,b)}.
\end{aligned}
\]